% chapters/ch07_can_bus.tex — Chapter 7: CAN Bus Deep Dive
% Scope (PLAN.md): frame fields, identifiers & arbitration, bit stuffing,
% error frames & states (active/passive/bus-off), ACK, bit timing & sample
% point, CAN FD awareness, higher layers (CANopen/J1939), drone servo control
% over CAN. Budget 45 pp, mix 28 Q&A / 13 code / 4 concept.
% All snippets verified gcc -Wall -Wextra -std=c11 (zero warnings).

\chapter{CAN Bus Deep Dive}
\label{ch:can-bus}

\begin{partnote}
\textbf{Scope} --- frame fields $\rightarrow$ identifiers \& bitwise
arbitration $\rightarrow$ bit stuffing $\rightarrow$ error frames \& fault
confinement (error-active / passive / bus-off) $\rightarrow$ the ACK mechanism
$\rightarrow$ bit timing \& sample point $\rightarrow$ CAN FD awareness
$\rightarrow$ higher layers (CANopen / J1939) $\rightarrow$ drone-servo control
over CAN. CAN is the candidate's lab bus (MKS drone CAN servos, CAN across the
UAV LRUs), so this is one of the strongest selling points --- it must be
bulletproof.
\end{partnote}

CAN earns its place in vehicles and aircraft because it is \emph{multi-master},
\emph{prioritised}, and \emph{robust}: any node can talk, the most urgent
message always wins the bus without collisions, and a five-layer error scheme
keeps one faulty node from taking down the bus. Every question below circles
that design.

%=====================================================================
\section{Primer}
%=====================================================================

\subsection{The bus and the bits}

CAN is a two-wire differential bus (\textbf{CAN\_H}, \textbf{CAN\_L}) with a
crucial asymmetry: a \textbf{dominant} bit (logical 0) wins over a
\textbf{recessive} bit (logical 1). Any node driving dominant pulls the bus
dominant --- wired-AND, like I2C but for arbitration. Nodes are not addressed;
instead each \emph{message} carries an \textbf{identifier} that serves both as
its priority and as a content label that receivers filter on. The bus is
terminated 120\,$\Omega$ at both ends (Chapter 5's transmission-line rule).

\subsection{Arbitration without collisions}

Because dominant beats recessive, CAN does collision-free arbitration: every
node that wants to transmit starts at the same time and sends its identifier
MSB-first while \emph{monitoring} the bus. A node sending a recessive (1) that
sees a dominant (0) has lost --- a higher-priority message is using the same bit
position --- so it backs off and becomes a receiver, \emph{without} any data
being destroyed. The winner (lowest numeric identifier = most dominant bits
early) continues seamlessly. This is the heart of CAN: priority is the ID, and
the lowest ID always wins, deterministically.

\subsection{Frame format, stuffing, CRC, ACK}

A standard data frame is: SOF (1 dominant) $\rightarrow$ 11-bit identifier
$\rightarrow$ RTR $\rightarrow$ control (IDE, r0, 4-bit DLC) $\rightarrow$
0--8 data bytes $\rightarrow$ 15-bit CRC + delimiter $\rightarrow$ ACK slot +
delimiter $\rightarrow$ 7-bit EOF, then the inter-frame space. Extended frames
use a 29-bit identifier. To keep receivers synchronised, the transmitter applies
\textbf{bit stuffing}: after five consecutive identical bits it inserts one of
opposite polarity (receivers remove it), guaranteeing enough edges for
resynchronisation. The \textbf{CRC} (15-bit) detects corruption; the
\textbf{ACK slot} is sent recessive by the transmitter and any node that
received the frame correctly pulls it dominant --- so the transmitter learns
\emph{at least one} node heard it.

\subsection{Bit timing, fault confinement, CAN FD}

Each bit is divided into \textbf{time quanta}: a sync segment, then phase
segments, with the \textbf{sample point} (typically ~75--87.5\%) where the bit
value is read --- placed late so signal propagation around the bus settles
first. \textbf{Fault confinement}: each node keeps a transmit and receive error
counter; on errors they climb and on success they fall. Error-\textbf{active}
(counter $<$128) signals errors actively; error-\textbf{passive}
($\geq$128) backs off; \textbf{bus-off} ($>$255 TEC) removes the node from the
bus entirely --- so a babbling node silences \emph{itself}, protecting the bus.
\textbf{CAN FD} (Flexible Data-rate) extends classic CAN with up to 64-byte
payloads and a faster data-phase bitrate, while keeping the same arbitration.

%=====================================================================
\section{Q\&A Bank}
%=====================================================================

\tierbanner{Tier 1 --- Screening (phone / HR-technical)}%
{Two to four sentences. CAN fundamentals --- expected cold for any
automotive/avionics role.}

\question{Q1.1}{What kind of bus is CAN?}
A multi-master, two-wire differential serial bus with message-based (not
node-addressed) communication and built-in priority arbitration and error
handling. Any node may transmit; the message identifier sets priority and is
what receivers filter on.

\question{Q1.2}{What are dominant and recessive bits?}
Dominant is logical 0 and ``wins'' --- if any node drives dominant, the bus is
dominant; recessive is logical 1 and only appears if every node lets the bus
float recessive. This wired-AND behaviour underlies arbitration and acks.

\question{Q1.3}{How does CAN decide who transmits when two nodes start
together?}
By bitwise arbitration on the identifier: each node sends its ID while watching
the bus, and a node that sends recessive but reads dominant has lost and backs
off. The lowest numeric ID (most dominant bits) wins, with no collision and no
data loss.

\question{Q1.4}{Is CAN addressed to nodes?}
No --- frames carry a message identifier, not a destination node address. Every
node hears every frame and uses acceptance filters to keep only the IDs it cares
about.

\question{Q1.5}{What is bit stuffing?}
After five consecutive bits of the same value, the transmitter inserts one
opposite bit (which receivers strip out). It guarantees regular edges so
receivers can stay synchronised on a bus with no separate clock.

\question{Q1.6}{How big is a CAN data payload?}
Classic CAN carries 0--8 data bytes per frame (DLC 0--8). CAN FD extends this to
up to 64 bytes.

\question{Q1.7}{What is the ACK slot?}
A bit the transmitter sends recessive; any node that received the frame without
error pulls it dominant. So the transmitter knows at least one node received it
correctly --- though not \emph{which} node.

\question{Q1.8}{Standard vs extended CAN identifier?}
Standard frames use an 11-bit identifier (2048 IDs); extended frames use 29 bits
(via the IDE bit), giving far more IDs --- used by higher layers like J1939.

\question{Q1.9}{Why is CAN good for vehicles and aircraft?}
Differential noise immunity, deterministic priority (the urgent message always
wins), multi-master flexibility, robust error detection and fault confinement,
and low wiring (one twisted pair for many nodes). It degrades gracefully ---
a faulty node removes itself.

\question{Q1.10}{What terminates a CAN bus and why?}
120\,$\Omega$ resistors at each physical end of the trunk, matching the cable
impedance to absorb reflections (Chapter 5). Two in parallel present ~60\,$\Omega$
to drivers; missing termination causes ringing and errors.

\tierbanner{Tier 2 --- Core technical (panel round)}%
{A short paragraph.}

\question{Q2.1}{Walk a concrete arbitration example.}
Suppose node A sends ID 0x100 and node B sends 0x200 simultaneously. They agree
on the early identical bits; at the first differing bit A sends 0 (dominant)
where B sends 1 (recessive). B reads the bus dominant while transmitting
recessive, so B loses arbitration and switches to receive; A continues without
even knowing a contest happened. 0x100 $<$ 0x200, so the lower ID won --- the
general rule.

\question{Q2.2}{Walk the fields of a standard CAN data frame in order.}
SOF (1 dominant start bit) $\rightarrow$ 11-bit identifier (MSB first)
$\rightarrow$ RTR (dominant for data, recessive for remote request)
$\rightarrow$ IDE (standard/extended) and r0 reserved $\rightarrow$ 4-bit DLC
(data length) $\rightarrow$ 0--8 data bytes $\rightarrow$ 15-bit CRC + recessive
delimiter $\rightarrow$ ACK slot + delimiter $\rightarrow$ 7 recessive EOF bits,
then $\geq$3 bits of inter-frame space. The receiver checks CRC, form, stuffing,
and acks in the ACK slot.

\question{Q2.3}{Explain CAN bit timing and the sample point.}
Each bit time is divided into time quanta: a fixed 1-tq sync segment, then
propagation + phase-segment-1 before the sample point, then phase-segment-2
after. The \textbf{sample point} is where the node reads the bit value, placed
late (~75--87.5\%) so the dominant/recessive level has propagated to the
furthest node and settled before sampling. You set it via the bit-rate
prescaler and TSEG1/TSEG2; both ends (the whole bus) must agree on bitrate, and
the sample point affects how much propagation delay and oscillator tolerance the
bus survives.

\question{Q2.4}{Compute a 500\,kbit/s configuration from a 40\,MHz CAN clock.}
Pick the prescaler and segments so
$\text{bitrate} = f_{\text{clk}} / ((\text{BRP}+1)\cdot N_{tq})$ with
$N_{tq}=1+\text{TSEG1}+\text{TSEG2}$. With BRP+1\,=\,5 and $N_{tq}=16$ (TSEG1=13,
TSEG2=2): $40\text{M}/(5\times16) = 500$\,kbit/s, and the sample point is
$(1+13)/16 = 87.5\%$. A late sample point like this is typical for longer buses
to tolerate propagation delay; you'd lower it if you needed more
oscillator-tolerance margin instead.

\question{Q2.5}{Explain the three fault-confinement states.}
Each node has a transmit error counter (TEC) and receive error counter (REC),
which rise on errors and fall on success. \textbf{Error-active} (counters
$<$128): normal --- the node signals errors with an active (dominant) error
flag. \textbf{Error-passive} (counter $\geq$128): the node still works but
signals errors passively and waits longer between transmits, so a flaky node
intrudes less. \textbf{Bus-off} (TEC $>$255): the node disconnects from the bus
entirely and only rejoins after a recovery sequence. The mechanism makes a
faulty node remove \emph{itself}, protecting the healthy traffic.

\question{Q2.6}{What is an error frame and what does it do?}
When a node detects an error (bit, stuff, CRC, form, or ACK error) it transmits
an error flag --- six dominant bits for an active node --- which violates bit
stuffing and forces every node to discard the current frame, and the
transmitter to retry. This globalises a locally-detected error so no node
accepts a corrupt frame, and it bumps the error counters that drive fault
confinement.

\question{Q2.7}{Name the five CAN error-detection mechanisms.}
(1) \textbf{Bit error}: a transmitter reads back a different bit than it sent
(outside arbitration/ack). (2) \textbf{Stuff error}: six consecutive identical
bits (stuffing violated). (3) \textbf{CRC error}: computed CRC $\neq$ received.
(4) \textbf{Form error}: a fixed-format field (delimiter, EOF) has the wrong
value. (5) \textbf{ACK error}: no node acknowledged. Together they give CAN its
very low undetected-error rate.

\question{Q2.8}{How do acceptance filters work and why do they matter?}
A receiver configures filter/mask register pairs: an incoming ID is accepted if
\code{(id \^{} filter) \& mask == 0} --- the mask selects which ID bits must
match. This lets a node ignore the bulk of bus traffic in hardware and interrupt
the CPU only for relevant messages, which is essential on a busy bus where
interrupting on every frame would swamp a small MCU. You design the ID scheme so
a single mask captures a group of related messages.

\question{Q2.9}{What is a remote frame (RTR)?}
A frame with the RTR bit recessive and no data, used to \emph{request} that the
node owning a given ID transmit its data frame. It's a pull mechanism, largely
deprecated in modern designs (and absent in CAN FD) because it complicates
timing and arbitration; most systems use periodic or event-driven data frames
instead.

\question{Q2.10}{What does CAN FD add, and what stays the same?}
CAN FD keeps the same arbitration, IDs, and bus, but adds a larger payload (up
to 64 bytes via a remapped DLC) and a \emph{bit-rate switch}: the data phase can
run faster than the arbitration phase, since arbitration only needs the slow
rate for bus-wide consensus while the payload --- sent by the single winner ---
can go fast. It also has a stronger CRC. The result is much higher throughput
without giving up CAN's deterministic arbitration.

\question{Q2.11}{How is a drone servo typically controlled over CAN?}
Via a higher-layer protocol (often CANopen-style or a vendor profile): the
flight controller broadcasts a command message (target position/speed) with a
known ID at a fixed rate, and the servo publishes status/feedback (actual
position, current, temperature, faults) on its own ID. Arbitration guarantees
the high-priority command and safety messages win the bus; acceptance filters
let each servo grab only its command. This is exactly the MKS drone CAN servo
interaction on the resume.

\question{Q2.12}{What is CANopen and what problem does it solve?}
CANopen is a standardised application layer over CAN that defines \emph{how} IDs
and data are organised: an Object Dictionary (addressable parameters), PDOs
(process data objects --- real-time data mapped to specific COB-IDs), SDOs
(service data for configuration), and NMT (network management/heartbeat). It
turns raw CAN frames into interoperable device profiles (e.g.\ for motors/
servos), so devices from different vendors cooperate without each team
inventing an ID map.

\tierbanner{Tier 3 --- Trap \& depth (principal-engineer level)}%
{A paragraph plus the trap. The defensible depth behind ``CAN servo lab work.''}

\question{Q3.1}{A node keeps going bus-off on a busy bus. Walk the root-cause
hunt.}
Bus-off means its TEC exceeded 255 --- it repeatedly detected errors on its own
transmissions. Causes, in order: (1) \textbf{bit-timing/bitrate mismatch} ---
even a small sample-point or bitrate disagreement makes this node see bit/stuff
errors while others are fine; (2) \textbf{termination} --- missing or wrong
120\,$\Omega$ causes reflections that corrupt edges, worst for the node
electrically furthest; (3) \textbf{a single faulty transceiver or stuck
driver}; (4) \textbf{ground offset / common-mode} beyond the transceiver range.
Method: scope CAN\_H/CAN\_L at the offending node for ringing (termination) and
check its bit-timing registers against the others, and confirm with a bus
analyser whether \emph{only} this node errors (its config) or everyone does (bus
fault). \trap the junior reflex is ``reset the node''; bus-off is a \emph{symptom
of physical or timing} disagreement, and the disciplined fix localises it on the
scope plus the bit-timing config --- the oscilloscope root-cause work on the
resume.

\question{Q3.2}{Why is the sample point a design trade, and what breaks if it's
wrong across the bus?}
The sample point must be late enough that a bit driven by the furthest node has
propagated across the whole bus (cable delay + transceiver delay) and settled
before everyone samples --- otherwise nodes read the bit at different effective
times and disagree. But placing it very late shrinks the margin for oscillator
tolerance (phase-segment-2 is also where resynchronisation jitter is absorbed).
So you trade bus length/propagation against clock-tolerance. If nodes are
configured with \emph{different} sample points or bitrates, they intermittently
disagree on marginal bits $\rightarrow$ stuff/bit errors $\rightarrow$ error
frames $\rightarrow$ one node climbing to bus-off. \trap the trap is assuming
``same bitrate'' is enough; the \emph{sample point} (TSEG1/TSEG2 split) must
also match across the bus, which is the subtle cause of a single node that
errors while the rest are happy.

\question{Q3.3}{Arbitration guarantees the lowest ID wins --- how can that
\emph{starve} a node, and how do you design around it?}
If a high-priority ID is transmitted continuously (or too frequently), it can
keep winning arbitration and a lower-priority node may rarely get the bus ---
priority is strict, so there's no fairness. The fix is at design time: budget the
bus load (sum of each periodic message's frame time $\times$ rate must leave
headroom), assign IDs by genuine real-time urgency (safety/control low,
diagnostics high), rate-limit chatty senders, and verify worst-case latency for
each message against its deadline. \trap the trap is treating ID assignment as
arbitrary labelling; on CAN the ID \emph{is} the scheduling priority, so a
careless map can starve a real function. Real avionics/automotive design does a
formal bus-load and worst-case-latency analysis --- the kind of integration
rigour the resume's V\&V work implies.

\question{Q3.4}{CAN's ACK only says ``someone heard it.'' Why is that
insufficient for a command to a specific servo, and what do you layer on top?}
A dominant ACK is pulled by \emph{any} receiving node, so a transmitter learns
the frame was well-formed and at least one node got it --- not that the
\emph{intended} servo received or acted on it. For a command you need
application-level confirmation: the servo publishes feedback (achieved position/
status) that the controller checks, plus heartbeats/node-guarding so a silent or
dead servo is detected, and a safe default (hold/neutral) on command timeout.
\trap the trap is trusting the CAN ACK as delivery-to-destination confirmation;
it isn't. The robust control loop closes on the servo's \emph{feedback} message
and treats missing feedback as a fault --- the same ``don't trust the
link-layer, validate at the application layer'' discipline as the SPI/UART
chapters, and the basis of safe actuator control (Chapter 10).

\question{Q3.5}{Estimate whether a 500\,kbit/s bus can carry 20 servos at
1\,kHz command + feedback, and what dominates.}
A standard 8-byte frame is ~111 bits without stuffing, ~135 worst-case with
stuffing, so at 500\,kbit/s a frame is ~222\,$\mu$s nominal (~270\,$\mu$s worst
case). 20 servos $\times$ (1 command + 1 feedback) at 1\,kHz = 40000 frames/s
$\times$ ~222\,$\mu$s $\approx$ 8.9 bus-seconds per second --- \emph{impossible}
(way over 100\% load). So either the rate, node count, or payload must drop, or
you move to CAN FD (64-byte frames batch many setpoints, and the fast data phase
slashes per-byte time) or a higher bitrate. \trap the trap is sizing by data
bytes alone and forgetting the ~44 bits of per-frame overhead plus stuffing and
inter-frame space, which dominate for small frequent messages --- exactly why
control buses either aggregate setpoints or move to CAN FD. Showing this
back-of-envelope load calc is what marks real integration judgement.

%=====================================================================
\section{Coding Gym}
%=====================================================================

\begin{partnote}
Thirteen host-compilable problems (verified under
\code{gcc -Wall -Wextra -std=c11}) covering the bit-level and protocol
arithmetic of CAN: arbitration, stuffing/de-stuffing, CRC-15, frame sizing and
bus load, bit timing, fault confinement, filtering, and higher-layer ID
decoding.
\end{partnote}

%---------------------------------------------------------------
\begin{gymproblem}{Arbitration winner}

\textbf{Problem.} Given two identifiers transmitted simultaneously, return
which wins.

\textbf{Hints.}
\begin{itemize}
\item Lowest numeric ID wins (most dominant bits early).
\end{itemize}

\attempt

\textbf{Solution.}
\begin{cgym}
#include <stdint.h>
#include <stdio.h>

static int can_wins(uint32_t a, uint32_t b) {
    return a < b;       /* lower ID = more dominant bits = higher priority */
}

int main(void) {
    printf("A wins=%d\n", can_wins(0x100u, 0x200u));   /* 1 */
    return 0;
}
\end{cgym}

The whole genius of CAN is that this comparison happens \emph{on the wire},
bit by bit, with no collision: the loser detects a dominant where it sent
recessive and silently becomes a receiver. So priority is simply the numeric ID
--- which is why the ID map \emph{is} your real-time schedule (Q3.3).

\followups
\begin{enumerate}
\item \emph{``Find the first differing bit where B loses.''} $\rightarrow$
  \code{clz}/scan of \code{a\^{}b} from the MSB.
\item \emph{``Two nodes with the same ID?''} $\rightarrow$ Forbidden --- they'd
  both ``win'' and their differing data causes bit errors; IDs must be unique
  per sender.
\item \emph{``Extended (29-bit) vs standard in arbitration?''} $\rightarrow$ A
  standard frame wins a tie at the RTR/IDE point against an extended one with
  the same 11-bit prefix.
\end{enumerate}
\end{gymproblem}

%---------------------------------------------------------------
\begin{gymproblem}{Bit stuffing}

\textbf{Problem.} Insert a stuff bit (opposite polarity) after every five
consecutive identical bits.

\textbf{Hints.}
\begin{itemize}
\item Track the current run length and last value; on the 5th, emit the
  opposite and reset the run to 1.
\end{itemize}

\attempt

\textbf{Solution.}
\begin{cgym}
#include <stdint.h>
#include <stdio.h>

static int stuff_bits(const uint8_t *in, int n, uint8_t *out) {
    int o = 0, run = 0, last = -1;
    for (int i = 0; i < n; i++) {
        int b = in[i];
        out[o++] = (uint8_t)b;
        if (b == last) run++; else { last = b; run = 1; }
        if (run == 5) { int s = !b; out[o++] = (uint8_t)s; last = s; run = 1; }
    }
    return o;
}

int main(void) {
    uint8_t in[6] = {1,1,1,1,1,1}, out[16];
    int n = stuff_bits(in, 6, out);
    printf("len=%d seq=", n);
    for (int i = 0; i < n; i++) printf("%d", out[i]);   /* 1111101 */
    printf("\n");
    return 0;
}
\end{cgym}

Six ones become \code{1111101}: after the fifth 1 a 0 is stuffed, then the sixth
1 follows. Stuffing guarantees an edge at least every six bits so receivers can
resynchronise their bit clock on a bus with no separate clock line --- and a
\emph{missing} stuff bit (six identical) is itself a detectable error.

\followups
\begin{enumerate}
\item \emph{``Why does worst-case framing grow with stuffing?''} $\rightarrow$
  Each run of 5 adds a bit; the worst case (problem 12) inflates the frame.
\item \emph{``Which fields are stuffed?''} $\rightarrow$ SOF through CRC; the
  delimiters, ACK, EOF are fixed-form (not stuffed).
\item \emph{``Six identical bits on the wire means?''} $\rightarrow$ A stuff
  error (or an error flag) --- an intentional stuffing violation.
\end{enumerate}
\end{gymproblem}

%---------------------------------------------------------------
\begin{gymproblem}{De-stuffing (round-trip)}

\textbf{Problem.} Remove stuff bits, recovering the original bit stream.

\textbf{Hints.}
\begin{itemize}
\item After five identical bits, the next bit is a stuff bit --- skip it.
\end{itemize}

\attempt

\textbf{Solution.}
\begin{cgym}
#include <stdint.h>
#include <stdio.h>

static int destuff_bits(const uint8_t *in, int n, uint8_t *out) {
    int o = 0, run = 0, last = -1, i = 0;
    while (i < n) {
        int b = in[i++];
        out[o++] = (uint8_t)b;
        if (b == last) run++; else { last = b; run = 1; }
        if (run == 5 && i < n) { last = in[i]; run = 1; i++; }   /* drop stuff bit */
    }
    return o;
}

int main(void) {
    uint8_t stuffed[7] = {1,1,1,1,1,0,1}, out[16];
    int n = destuff_bits(stuffed, 7, out);
    printf("len=%d seq=", n);
    for (int i = 0; i < n; i++) printf("%d", out[i]);   /* 111111 */
    printf("\n");
    return 0;
}
\end{cgym}

De-stuffing is the exact inverse: count five identical received bits, then
discard the following (stuff) bit. A correct round-trip (\code{stuff} then
\code{destuff} returns the original, verified) is the property a receiver relies
on; if the de-stuffer ever finds a \emph{sixth} identical bit instead of the
expected opposite, that's a stuff error and the frame is rejected.

\followups
\begin{enumerate}
\item \emph{``What if the stuff bit was corrupted to match?''} $\rightarrow$
  Then six identical appear $\rightarrow$ stuff error detected; CRC also guards.
\item \emph{``Why count the stuffed bit's value as the new run start?''}
  $\rightarrow$ It can itself begin a new run that may need stuffing.
\item \emph{``De-stuff in hardware vs software?''} $\rightarrow$ The CAN
  controller does it; this models what it performs.
\end{enumerate}
\end{gymproblem}

%---------------------------------------------------------------
\begin{gymproblem}{CAN CRC-15}

\textbf{Problem.} Compute the CAN 15-bit CRC (polynomial 0x4599) over a bit
sequence.

\textbf{Hints.}
\begin{itemize}
\item Classic shift-register CRC: XOR the incoming bit with the top CRC bit;
  shift; conditionally XOR the polynomial.
\end{itemize}

\attempt

\textbf{Solution.}
\begin{cgym}
#include <stdint.h>
#include <stdio.h>

static uint16_t can_crc15(const uint8_t *bits, int n) {
    uint16_t crc = 0;
    for (int i = 0; i < n; i++) {
        uint16_t inb = (uint16_t)(((crc >> 14) & 1u) ^ (bits[i] & 1u));
        crc = (uint16_t)((crc << 1) & 0x7FFFu);    /* 15-bit register */
        if (inb) crc ^= 0x4599u;                   /* CAN polynomial */
    }
    return crc;
}

int main(void) {
    uint8_t bits[10] = {1,0,1,1,0,0,1,0,1,1};
    printf("crc=0x%04X\n", can_crc15(bits, 10));   /* 0x6336 */
    return 0;
}
\end{cgym}

The 15-bit CRC over the frame (SOF through the data field) is what gives CAN its
famously low undetected-error rate --- it catches all single- and double-bit
errors and all burst errors up to 15 bits. The transmitter appends it; every
receiver recomputes and a mismatch triggers an error frame and a retransmission.
This is the link-layer integrity beneath any higher-layer protocol.

\followups
\begin{enumerate}
\item \emph{``Why a CRC and not a checksum?''} $\rightarrow$ CRC catches burst
  errors a sum misses --- critical on a noisy vehicle bus.
\item \emph{``CAN FD CRC?''} $\rightarrow$ Longer (17- or 21-bit) for the larger
  payloads, plus stuff-bit counting.
\item \emph{``Does the CRC cover stuff bits?''} $\rightarrow$ No --- it's
  computed on the de-stuffed bit stream.
\end{enumerate}
\end{gymproblem}

%---------------------------------------------------------------
\begin{gymproblem}{Standard frame bit length}

\textbf{Problem.} Compute the number of bits in a standard CAN frame (no
stuffing) for a given DLC.

\textbf{Hints.}
\begin{itemize}
\item Sum the fixed fields + 8$\times$DLC data bits + EOF + IFS.
\end{itemize}

\attempt

\textbf{Solution.}
\begin{cgym}
#include <stdio.h>

static int can_std_frame_bits(int dlc) {
    return 1 + 11 + 1 + 1 + 1 + 4   /* SOF, ID, RTR, IDE, r0, DLC */
         + 8 * dlc                   /* data */
         + 15 + 1 + 1 + 1 + 7 + 3;   /* CRC, CRCdelim, ACK, ACKdelim, EOF, IFS */
}

int main(void) {
    printf("bits(dlc=8)=%d\n", can_std_frame_bits(8));   /* 111 */
    return 0;
}
\end{cgym}

An 8-byte standard frame is ~111 bits, of which ~44 are \emph{overhead}
(everything but the data). That overhead is the key number for bus-load
estimation: for small frequent messages the fixed cost dominates, which is why
CAN FD's 64-byte payload (amortising the overhead) and faster data phase win for
high-throughput control (Q3.5).

\followups
\begin{enumerate}
\item \emph{``Extended frame?''} $\rightarrow$ Add 18 ID bits + SRR + IDE
  changes ($\approx$20 more bits).
\item \emph{``Why does overhead dominate small frames?''} $\rightarrow$ ~44
  fixed bits vs 8--64 data bits --- the ratio is terrible for tiny payloads.
\item \emph{``Where does the worst-case 135 come from?''} $\rightarrow$ Maximum
  bit-stuffing on the stuffable fields (problem 12).
\end{enumerate}
\end{gymproblem}

%---------------------------------------------------------------
\begin{gymproblem}{Frame time and bus load}

\textbf{Problem.} Compute a frame's transmission time at a given bitrate.

\textbf{Hints.}
\begin{itemize}
\item time = bits / bitrate.
\end{itemize}

\attempt

\textbf{Solution.}
\begin{cgym}
#include <stdint.h>
#include <stdio.h>

static uint32_t can_frame_us(int frame_bits, uint32_t bitrate) {
    return (uint32_t)((uint64_t)frame_bits * 1000000u / bitrate);
}

int main(void) {
    printf("us=%u\n", can_frame_us(111, 500000u));   /* 222 us */
    return 0;
}
\end{cgym}

111 bits at 500\,kbit/s is 222\,$\mu$s. Bus load is then
$\sum(\text{frame\_us} \times \text{rate})$ over all periodic messages; keep it
well under 100\% (often $<$70--80\% for latency headroom). This single
calculation decides whether a proposed message set is even feasible --- the
back-of-envelope every CAN integrator runs before committing an ID map
(Q3.5).

\followups
\begin{enumerate}
\item \emph{``Add 30\% for worst-case stuffing --- new time?''} $\rightarrow$
  ~289\,$\mu$s; use the worst case for deadline analysis.
\item \emph{``Why keep load under ~80\%?''} $\rightarrow$ Headroom bounds
  worst-case arbitration latency for low-priority messages.
\item \emph{``Highest-priority message's worst-case latency?''} $\rightarrow$
  At most one in-progress frame ahead of it (non-preemptive), then it wins.
\end{enumerate}
\end{gymproblem}

%---------------------------------------------------------------
\begin{gymproblem}{Bit timing: bitrate and sample point}

\textbf{Problem.} From the CAN clock, prescaler, and segment counts, compute the
bitrate and the sample-point percentage.

\textbf{Hints.}
\begin{itemize}
\item $N_{tq}=1+\text{TSEG1}+\text{TSEG2}$; bitrate = clk / ((BRP+1)$\times
  N_{tq}$); sample = (1+TSEG1)/$N_{tq}$.
\end{itemize}

\attempt

\textbf{Solution.}
\begin{cgym}
#include <stdint.h>
#include <stdio.h>

static uint32_t can_bitrate(uint32_t fclk, int brp, int tseg1, int tseg2) {
    int tq = 1 + tseg1 + tseg2;
    return fclk / ((uint32_t)(brp + 1) * (uint32_t)tq);
}
static int can_sample_pct(int tseg1, int tseg2) {
    int tq = 1 + tseg1 + tseg2;
    return (1 + tseg1) * 100 / tq;
}

int main(void) {
    printf("bitrate=%u sample=%d%%\n",
           can_bitrate(40000000u, 4, 13, 2), can_sample_pct(13, 2));
    /* 500000 bits/s, 87% sample point */
    return 0;
}
\end{cgym}

40\,MHz with BRP+1=5 and 16 time-quanta (TSEG1=13, TSEG2=2) gives exactly
500\,kbit/s at an 87.5\% sample point. Every node on the bus must match the
bitrate \emph{and} keep a compatible sample point (Q3.2) --- mismatch is the
quiet cause of one node erroring its way to bus-off while the rest run fine.

\followups
\begin{enumerate}
\item \emph{``Move the sample point earlier --- why would you?''} $\rightarrow$
  More margin for oscillator tolerance, less for propagation; a short bus can
  afford it.
\item \emph{``Same bitrate, different BRP/segments --- compatible?''}
  $\rightarrow$ Bitrate matches but sample points may differ; align them.
\item \emph{``SJW (resync jump width)?''} $\rightarrow$ How many tq the node may
  shrink/grow a segment to resync; bounds tolerable clock drift.
\end{enumerate}
\end{gymproblem}

%---------------------------------------------------------------
\begin{gymproblem}{Fault-confinement state}

\textbf{Problem.} Map a transmit error counter to error-active, error-passive,
or bus-off.

\textbf{Hints.}
\begin{itemize}
\item $<$128 active, $\geq$128 passive, $>$255 bus-off.
\end{itemize}

\attempt

\textbf{Solution.}
\begin{cgym}
#include <stdio.h>

typedef enum { ERR_ACTIVE, ERR_PASSIVE, BUS_OFF } can_err_t;

static can_err_t can_state(int tec) {
    if (tec > 255) return BUS_OFF;
    if (tec > 127) return ERR_PASSIVE;
    return ERR_ACTIVE;
}

int main(void) {
    printf("%d %d %d\n", can_state(96), can_state(200), can_state(300));
    /* ERR_ACTIVE(0) ERR_PASSIVE(1) BUS_OFF(2) */
    return 0;
}
\end{cgym}

This is CAN's self-protection: a node that keeps erroring climbs from active to
passive (intrudes less on the bus) to bus-off (disconnects), so a babbling or
broken node silences itself rather than jamming healthy traffic. Counters fall on
successful frames, so a transient glitch recovers; only persistent faults reach
bus-off, which requires an explicit recovery sequence to rejoin.

\followups
\begin{enumerate}
\item \emph{``How does a node leave bus-off?''} $\rightarrow$ After 128
  occurrences of 11 recessive bits (bus idle), it may reset its counters and
  rejoin.
\item \emph{``Active vs passive error flag difference?''} $\rightarrow$ Active =
  6 dominant bits (forces all nodes); passive = 6 recessive (doesn't disturb a
  healthy bus).
\item \emph{``Why count TX errors more heavily?''} $\rightarrow$ A node erroring
  its \emph{own} transmissions is more likely the faulty one.
\end{enumerate}
\end{gymproblem}

%---------------------------------------------------------------
\begin{gymproblem}{CANopen COB-ID decode}

\textbf{Problem.} Split a CANopen COB-ID into its function code and node ID.

\textbf{Hints.}
\begin{itemize}
\item Node ID = low 7 bits; function code = next 4 bits.
\end{itemize}

\attempt

\textbf{Solution.}
\begin{cgym}
#include <stdint.h>
#include <stdio.h>

static void cob_split(uint16_t cob, int *func, int *node) {
    *func = (cob >> 7) & 0xF;     /* function code */
    *node = cob & 0x7F;           /* node ID (1..127) */
}

int main(void) {
    int f, n;
    cob_split(0x181, &f, &n);
    printf("func=%d node=%d\n", f, n);   /* func=3 (TPDO1), node=1 */
    return 0;
}
\end{cgym}

CANopen layers structure onto the raw 11-bit ID: the top 4 bits are a function
code (0x3 = Transmit-PDO1, 0x4 = Receive-PDO1, etc.) and the low 7 are the node
ID. So 0x181 is node 1's first transmit-PDO --- a servo broadcasting its
feedback, for instance. This predefined connection set is why CANopen devices
interoperate without a custom ID map (Q2.12).

\followups
\begin{enumerate}
\item \emph{``What COB-ID is node 5's RPDO1?''} $\rightarrow$ Function 0x4
  $\rightarrow$ \code{(0x4<<7)|5 = 0x205}.
\item \emph{``Node 0 --- valid?''} $\rightarrow$ No; node IDs are 1--127, 0 is
  reserved (broadcast/NMT).
\item \emph{``What's at function code 0x0?''} $\rightarrow$ NMT / sync ---
  network management, highest priority.
\end{enumerate}
\end{gymproblem}

%---------------------------------------------------------------
\begin{gymproblem}{Acceptance filter match}

\textbf{Problem.} Decide whether a received ID passes a filter/mask pair.

\textbf{Hints.}
\begin{itemize}
\item Pass when \code{((id \^{} filter) \& mask) == 0}.
\end{itemize}

\attempt

\textbf{Solution.}
\begin{cgym}
#include <stdint.h>
#include <stdio.h>

static int filt_match(uint32_t id, uint32_t filter, uint32_t mask) {
    return ((id ^ filter) & mask) == 0;   /* masked bits must equal the filter */
}

int main(void) {
    printf("%d %d\n",
           filt_match(0x123, 0x120, 0x7F0),   /* low nibble don't-care -> 1 */
           filt_match(0x133, 0x120, 0x7F0));  /* bit 4 differs -> 0 */
    return 0;
}
\end{cgym}

The mask selects which ID bits must match the filter; \code{0} bits in the mask
are ``don't care.'' Here mask \code{0x7F0} ignores the low nibble, so \code{0x120}
--\code{0x12F} all pass --- one filter capturing a group of related messages.
Hardware applies this before raising an interrupt, so the CPU only wakes for
relevant traffic --- essential on a busy bus (Q2.8). Designing the ID layout so a
single mask grabs each functional group is part of the bus design.

\followups
\begin{enumerate}
\item \emph{``Match exactly one ID?''} $\rightarrow$ Set mask to all-ones over
  the ID width.
\item \emph{``Accept everything?''} $\rightarrow$ Mask = 0 (all don't-care).
\item \emph{``Why filter in hardware?''} $\rightarrow$ To avoid an interrupt per
  frame on a high-rate bus; the CPU can't keep up otherwise.
\end{enumerate}
\end{gymproblem}

%---------------------------------------------------------------
\begin{gymproblem}{DLC to data length (CAN FD mapping)}

\textbf{Problem.} Convert a 4-bit DLC to the number of data bytes, including the
CAN FD extended lengths.

\textbf{Hints.}
\begin{itemize}
\item 0--8 map directly; 9--15 map to 12,16,20,24,32,48,64.
\end{itemize}

\attempt

\textbf{Solution.}
\begin{cgym}
#include <stdio.h>

static int dlc_to_len(int dlc) {
    static const int m[16] = {0,1,2,3,4,5,6,7,8,12,16,20,24,32,48,64};
    return (dlc >= 0 && dlc < 16) ? m[dlc] : 0;
}

int main(void) {
    printf("dlc9=%d dlc15=%d\n", dlc_to_len(9), dlc_to_len(15));   /* 12 64 */
    return 0;
}
\end{cgym}

In classic CAN a DLC above 8 still means 8 bytes; CAN FD repurposes codes 9--15
for the larger payloads (12 up to 64 bytes). A receiver must use this table, not
``DLC = byte count,'' or it mis-parses FD frames. The big jumps (24$\rightarrow$
32$\rightarrow$48$\rightarrow$64) are why FD batches many small setpoints into
one frame to beat the per-frame overhead.

\followups
\begin{enumerate}
\item \emph{``Classic CAN with DLC=12 --- how many bytes?''} $\rightarrow$ 8
  (values $>$8 clamp to 8 in classic CAN).
\item \emph{``Why non-contiguous lengths?''} $\rightarrow$ A 4-bit field can't
  encode every value to 64, so FD picks useful steps.
\item \emph{``Pack 6 servo setpoints (2 bytes each) into one frame.''}
  $\rightarrow$ 12 bytes $\rightarrow$ DLC 9; one FD frame instead of six classic
  frames.
\end{enumerate}
\end{gymproblem}

%---------------------------------------------------------------
\begin{gymproblem}{Worst-case frame length with stuffing}

\textbf{Problem.} Given the number of stuffable bits, compute the worst-case
length after maximum bit stuffing.

\textbf{Hints.}
\begin{itemize}
\item Worst case adds one stuff bit per four bits after the first:
  \code{n + (n-1)/4}.
\end{itemize}

\attempt

\textbf{Solution.}
\begin{cgym}
#include <stdio.h>

static int worst_case_stuffed(int stuffable_bits) {
    return stuffable_bits + (stuffable_bits - 1) / 4;
}

int main(void) {
    printf("worst(34)=%d\n", worst_case_stuffed(34));   /* 42 */
    return 0;
}
\end{cgym}

Maximum stuffing inserts a bit roughly every four (a run of 5 forces one, then
the pattern can repeat), so the stuffable region grows by \code{(n-1)/4}. Applied
to a full 8-byte frame's stuffable fields this is how the ~111-bit nominal frame
reaches the famous ~135-bit worst case --- the number you \emph{must} use for
real-time latency analysis, because an adversarial data pattern can always force
it.

\followups
\begin{enumerate}
\item \emph{``Why use worst case, not average, for deadlines?''} $\rightarrow$
  Real-time guarantees must hold for the worst data pattern, not the typical
  one.
\item \emph{``Can you avoid stuffing?''} $\rightarrow$ No --- it's mandatory;
  you can only bound its worst case.
\item \emph{``Does the fixed-form EOF/ACK get stuffed?''} $\rightarrow$ No, so
  they're excluded from the stuffable count.
\end{enumerate}
\end{gymproblem}

%---------------------------------------------------------------
\begin{gymproblem}{J1939 PGN extraction}

\textbf{Problem.} Extract the Parameter Group Number from a 29-bit J1939
identifier.

\textbf{Hints.}
\begin{itemize}
\item PF (PDU format) = bits 16--23. If PF $<$ 240 (PDU1), PS is a destination
  address (not in the PGN); if PF $\geq$ 240 (PDU2), PS is the group extension.
\end{itemize}

\attempt

\textbf{Solution.}
\begin{cgym}
#include <stdint.h>
#include <stdio.h>

static uint32_t j1939_pgn(uint32_t ext_id) {
    uint32_t pf = (ext_id >> 16) & 0xFFu;
    uint32_t ps = (ext_id >> 8)  & 0xFFu;
    uint32_t dp = (ext_id >> 24) & 0x3u;     /* data page bits */
    return (pf < 240u) ? ((dp << 16) | (pf << 8))          /* PDU1: dest addr excluded */
                       : ((dp << 16) | (pf << 8) | ps);    /* PDU2: PS is group ext   */
}

int main(void) {
    printf("pgn=%u\n", j1939_pgn(0x18FEF100u));   /* 65265 (vehicle speed etc.) */
    return 0;
}
\end{cgym}

J1939 (trucks, gensets, some UAV/marine systems) layers a structured 29-bit ID:
priority, data page, PF, PS, and source address. The PGN names the message
\emph{content} (e.g.\ 65265 is a standard vehicle-dynamics group). The
PDU1/PDU2 split --- whether PS is a destination address or part of the PGN ---
is the subtlety that trips up naive parsers. Recognising these higher layers
(CANopen, J1939) signals you've worked above raw frames.

\followups
\begin{enumerate}
\item \emph{``Extract the source address.''} $\rightarrow$ Low 8 bits of the ID.
\item \emph{``Priority?''} $\rightarrow$ Top 3 bits (bits 26--28).
\item \emph{``Why does PDU1 exclude PS from the PGN?''} $\rightarrow$ PS is the
  \emph{destination address} for point-to-point messages, not part of the group
  identity.
\end{enumerate}
\end{gymproblem}

%=====================================================================
\section{Link to Her Resume}
%=====================================================================

\begin{resumelink}
\textbf{Drone-servo CAN lab work:} ``I work with CAN servos and CAN across the
UAV LRUs. A servo takes a periodic command on a known ID and publishes feedback
on its own ID; arbitration guarantees the high-priority command and safety
messages win the bus, and acceptance filters let each servo grab only its
command in hardware.''

\medskip
\textbf{Bus integrity \& bring-up:} ``A node going bus-off is a timing or
physical disagreement, not a software glitch --- I scope CAN\_H/CAN\_L for
ringing (termination) and check the bit-timing and sample point against the rest
of the bus, the same oscilloscope root-cause method I use across protocols.''

\medskip
\textbf{Integration judgement:} ``Before committing an ID map I run a bus-load
and worst-case-latency check --- per-frame overhead and stuffing dominate small
frequent messages, so control buses aggregate setpoints or move to CAN FD.''

\medskip
\small Maps to the resume's CAN servo / multi-LRU work and oscilloscope-based
root-cause analysis. Keep claims to what the resume states; deep project
scripting is Chapter 15, from \code{inputs/INTAKE.md} only.
\end{resumelink}

%=====================================================================
\section{Self-Test Scorecard}
%=====================================================================

\begin{scorecard}
\begin{enumerate}
\item What makes a bit dominant, and why does that enable collision-free
  arbitration?
\item Which CAN ID wins arbitration, and what does the loser do?
\item What is bit stuffing and why is it needed?
\item Name the five CAN error-detection mechanisms.
\item What are the three fault-confinement states and their thresholds?
\item What does the ACK slot actually tell the transmitter --- and what does it
  not?
\item Compute the bitrate and sample point for 40\,MHz, BRP+1=5, TSEG1=13,
  TSEG2=2.
\item How does an acceptance filter/mask decide to accept an ID?
\item What does CAN FD add while keeping arbitration the same?
\item Roughly how many bits is an 8-byte standard frame, and why does that
  matter?
\end{enumerate}
\end{scorecard}

\begin{answerkey}
\begin{enumerate}
\item Dominant (0) overrides recessive (1) on the wired-AND bus, so a node that
  sends recessive but reads dominant has lost --- contention resolves bit by bit
  with no collision.
\item The lowest numeric ID wins; the loser stops transmitting and becomes a
  receiver, with no data lost.
\item Inserting an opposite bit after five identical ones; it guarantees edges
  for receiver resynchronisation on a clockless bus.
\item Bit error, stuff error, CRC error, form error, ACK error.
\item Error-active ($<$128), error-passive ($\geq$128), bus-off (TEC $>$255).
\item That the frame was well-formed and at least one node received it; not
  \emph{which} node, nor that the intended node acted on it.
\item 500\,kbit/s, 87.5\% sample point ($N_{tq}=16$, $(1+13)/16$).
\item Accept when \code{((id \^{} filter) \& mask) == 0} --- masked bits must
  match the filter, zero-mask bits are don't-care.
\item Up to 64-byte payloads and a faster data-phase bitrate (and stronger CRC),
  with the same arbitration/IDs.
\item ~111 bits (~44 of overhead); it sets bus load and shows why small frequent
  messages are overhead-dominated (push toward CAN FD/aggregation).
\end{enumerate}
\end{answerkey}

\begin{partnote}
\emph{End of Chapter 7.} Next: \textbf{Avionics Data Buses} ---
MIL-STD-1553B, ARINC 429, Ethernet for telemetry, and cross-channel/redundancy
concepts (CCDL/SPIL-type links), connecting CAN's ideas to the
flight-grade buses on the resume.
\end{partnote}
